\section{The all-degree inverse}
\label{sec:inverse}
For the rest of the algebra keep $g,\kappa_0,\kappa_1$ fixed and put
\begin{equation}\label{eq:constants}
\begin{gathered}
 L_b=\frac{g}{2c_bz},\qquad \mathsf R=1+2z,\qquad V_m=1+2mz,\\
 k_m=\frac4{2^m(m+1)(m!)^2},\qquad
 C_m=\frac{2^{2-m}}{(m+1)(m+2)(m!)^2}.
\end{gathered}
\end{equation}
Here $\mathsf R$ is not the raw first moment $R_b$.
Use column vectors $Q_r=(q_{0,r},q_{1,r})^t$,
$\xi_k=(\xi_{0,k},\xi_{1,k})^t$, and
$\eta_k=(\eta_{0,k},\eta_{1,k})^t$.

\begin{proposition}[Successive even and odd blocks]\label{prop:blocks}
For every $m\ge2$, $\xi_{m-1}$ depends only on $Q_3,\ldots,Q_{2m}$,
is affine in $Q_{2m}$, and
\begin{equation}\label{eq:evenblock}
 D_{Q_{2m}}\xi_{m-1}
 =-\begin{pmatrix}\mathsf R^m&V_m\\V_m&\mathsf R^m\end{pmatrix}
                      \diag(k_mL_0^m,k_mL_1^m).
\end{equation}
For every $m\ge1$, $\eta_m$ depends only on $Q_3,\ldots,Q_{2m+1}$,
is affine in $Q_{2m+1}$, and
\begin{equation}\label{eq:oddblock}
 D_{Q_{2m+1}}\eta_m
 =-\begin{pmatrix}c_1\mathsf R^m&V_m\\
                   V_m&c_0\mathsf R^m\end{pmatrix}
              \diag(2c_0C_mL_0^{m+1},2c_1C_mL_1^{m+1}).
\end{equation}
No evenness assumption is needed for either assertion. The symmetric
matrix in each formula is positive definite. In particular,
\begin{align}
 \mathsf R^m-V_m&=\sum_{r=2}^m\binom mr(2z)^r>0
                          &&(m\ge2),\label{eq:evenpositive}\\
 c_0c_1\mathsf R^{2m}-V_m^2
   &\ge z(1+2mz)^2>0 &&(m\ge1).\label{eq:oddpositive}
\end{align}
\end{proposition}

\subsection{Finite-jet dependence}
Vary one graph derivative of order $r\ge3$, keeping all lower graph
jets fixed. Stationarity of the intermediate point gives the leading
variation of the reduced action without differentiating that point.
With $w_0=(u+v)/(2c_o)$ one obtains
\begin{align}
 \partial_{q_{b,r}}E_b&=\frac{u^r+v^r}{r!}+O(|y|^{r+1}),
       &\partial_{q_{o,r}}E_b&=\frac{2w_0^r}{r!}+O(|y|^{r+1}),
                                                \label{eq:variationE}\\
 \partial_{q_{b,r}}b_b&=O(|y|^{r-1}),
       &\partial_{q_{o,r}}b_b&=-\frac{g}{c_o}
                    \frac{w_0^{r-2}}{(r-2)!}+O(|y|^{r-1}).
                                                \label{eq:variationb}
\end{align}
Indeed the graph perturbation changes the longitudinal flight length
to first order with coefficient one at the normal orbit; each endpoint
is counted once and the intermediate contact twice. The correction to
$w_0$ has order two, so it changes no displayed homogeneous term.
Twice differentiating the action and using $d_b^0=(2gc_o)^{-1}$
proves the twist formulas. This is where the contribution of the
opposite contact to the density enters.

Here and below the remainder statements mean Taylor remainders with
the finite number of derivatives needed at the chosen order. The
coefficient assertions can be justified on the fixed Morse disk in
Lemma~\ref{lem:flux}. Equivalently, at the scaling $y=hY$, $d=h^2$,
a graph jet of degree $r$ first enters the rescaled action and twist
with weight $h^{r-2}$. The mark adds a factor $h$. Every nonquadratic
unvaried correction adds at least one further power of $h$. Odd powers
integrate to zero after the fixed-domain construction.
Consequently $[d^{m-1}]G_b$ uses only graph jets through degree $2m$,
whereas $[d^m]J_b$ uses only those through degree $2m+1$.
At their respective highest degrees the displayed homogeneous
variations are integrated against the quadratic action and constant
twist. Products with any lower nonlinear correction have higher weight.
The resulting highest-jet derivative is independent of all jets of
order at least three. This proves affinity, not just linearization at
a symmetric reference contact.

To see explicitly that these are analytic finite-jet maps, differentiate
the stationary equation at the origin degree by degree. Its highest
unknown derivative is always multiplied by $2c_o/g>0$. The remaining
operations are finite sums, products, and division by positive
leading-geometry quantities. Integration of the finitely many terms on
the quadratic disk preserves analytic dependence. Neither analyticity
nor evenness of the original smooth graph is needed for this argument.

\subsection{Quadratic moments}
For a linear form $\ell(y)$ set
$\nu(\ell)=\ell^tH_b^{-1}\ell$, and for two forms let
$\operatorname{cov}(t,\ell)=t^tH_b^{-1}\ell$.
Inverting \eqref{eq:H} gives
\begin{equation}\label{eq:covariances}
 H_b^{-1}=L_b\begin{pmatrix}\mathsf R&1\\1&\mathsf R\end{pmatrix},
 \qquad \nu(u)=L_b\mathsf R,\quad \nu(w_0)=L_o.
\end{equation}
With $t=u+v$,
\begin{equation}\label{eq:tcov}
 \operatorname{cov}(t,u)=\operatorname{cov}(t,v)
        =2c_bc_oL_b,\qquad t=2c_ow_0.
\end{equation}
Rotation after the transformation $y=H_b^{-1/2}x$ gives
\begin{align}
 \frac1{I_0(d)}\int_{E_{b,0}<d}\ell^{2r}\dd y
 &=\frac{2(2r)!\nu(\ell)^r}{2^r(r!)^2(r+1)}d^{r-1},
                                          \label{eq:diskmoment}\\
 \frac1{I_0(d)}\int(d-E_{b,0})_+\ell^{2r}\dd y
 &=\frac{2\binom{2r}{r}\nu(\ell)^r}{2^r(r+1)(r+2)}d^r.
                                          \label{eq:weightedmoment}
\end{align}
For example, the angular mean is
$\binom{2r}{r}/4^r$ and the unweighted radial integral is
$(2d)^{r+1}/(2r+2)$. For either radial weight and every $r\ge0$,
\begin{equation}\label{eq:crossmoment}
 \int t\ell^{2r+1}\,\mathrm{weight}
   =\frac{\operatorname{cov}(t,\ell)}{\nu(\ell)}
             \int\ell^{2r+2}\,\mathrm{weight}.
\end{equation}
To prove this identity, decompose $t$ in the transformed coordinates
into its component along $\ell$ and the orthogonal component. The latter
integrates to zero by reflection. All these formulas are exact.

\subsection{Computation of the blocks}
The first action variation in either integral is the negative integral
of the action perturbation over $E_{b,0}<d$, with the appropriate mark.
There is no extra moving-boundary term because the residual weight
vanishes on the boundary.
For $r=2m$, the two endpoint powers in \eqref{eq:variationE} give
$-k_mL_b^m\mathsf R^m$. The opposite-contact action gives
$-k_mL_o^m$. Its twist variation, by \eqref{eq:weightedmoment}, gives
$-k_mL_o^m(2mz)$, since $g/(c_oL_o)=2z$.
The sum is the first row of \eqref{eq:evenblock}; reversing the
contacts gives the other row. This calculation remains valid when
lower odd jets are nonzero, by the weight argument above.

For $r=2m+1$, use the mark $t=u+v$. Equations
\eqref{eq:diskmoment} and \eqref{eq:crossmoment}, with
$r$ there replaced by $m+1$, yield the own-contact action contribution
\begin{equation}\label{eq:ownodd}
 -C_m\operatorname{cov}(t,u)\nu(u)^m
       =-2c_bc_oC_mL_b^{m+1}\mathsf R^m.
\end{equation}
The opposite-contact action contributes $-2c_oC_mL_o^{m+1}$.
For its twist, use $t=2c_ow_0$ and \eqref{eq:weightedmoment}:
\begin{align*}
 &-\frac{g}{c_o(2m-1)!I_0(d)}
       \int t\,w_0^{2m-1}(d-E_{b,0})_+\dd y\\
 &\hspace{1cm}=-2m g C_mL_o^m d^m
       =-4mc_ozC_mL_o^{m+1}d^m.
\end{align*}
Thus the opposite-contact coefficient is
$-2c_oC_mL_o^{m+1}(1+2mz)$. Together with \eqref{eq:ownodd}
this proves \eqref{eq:oddblock}.

For the even symmetric matrix, \eqref{eq:evenpositive} and the
positive sum of the two entries prove positivity. For the odd one,
Bernoulli's inequality gives $\mathsf R^m\ge1+2mz$ for all
$m\ge1$. Since $c_0c_1=1+z$, its determinant is at least
$z(1+2mz)^2$. Its diagonal entries are positive. This proves
Proposition~\ref{prop:blocks} at every order.

At $g=1$, $c_0=c_1=2$, the first odd block is
\begin{equation}\label{eq:cubicexample}
 D_{Q_3}\eta_1=-\frac7{108}
                  \begin{pmatrix}2&1\\1&2\end{pmatrix}.
\end{equation}
Its nonzero antisymmetric eigenvalue is already visible at cubic
order. The first even block is
$D_{Q_4}\xi_1=-\frac1{1728}\left(\begin{smallmatrix}49&13\\13&49\end{smallmatrix}\right)$.
These examples are normalizations of the proof, not a substitute for it.

\subsection{Recursive recovery and analytic determination}
Order the data by graph degree as
\[
 \eta_1,\ \xi_1,\ \eta_2,\ \xi_2,\ldots .
\]
At degree $r$, let $Y_r$ be the corresponding two-component datum,
let $B_r$ be its block in Proposition~\ref{prop:blocks}, and let
$F_r(Q_3,\ldots,Q_{r-1})$ be its value when $Q_r=0$.
Then the inverse is the exact recursion
\begin{equation}\label{eq:recursion}
 Q_r=B_r^{-1}\{Y_r-F_r(Q_3,\ldots,Q_{r-1})\},\qquad r\ge3.
\end{equation}
Every inverse block exists. A fixed finite collection of its denominators
has a positive minimum on a compact positive leading-geometry box.
The recursion is analytic; on sufficiently small compact coordinate
neighborhoods the map and inverse have bounded derivatives. Integration
along segments gives the local bi-Lipschitz conclusion. This proves
part (ii) of Theorem~\ref{thm:main}.

For analytic contacts, the four raw germs first determine $A$ and then
all the coefficients in \eqref{eq:coefficients}. Recursion
\eqref{eq:recursion} identifies every graph derivative. Analyticity
identifies the graph germs. To obtain the boundary images, continue the
common analytic regular curve along arclength. At each point the local
graph identity continues uniquely; compactness and connectedness of
the closed embedded boundaries continue it around the entire curve.
This proves part (iii). The continuation concerns the participating
boundaries, not an obstacle never visited by the selected channel.

\begin{remark}[Why the signed information is essential]
The transformation $\psi_b(y)\mapsto\psi_b(-y)$ at both contacts
sends $(u,w,v)$ to $(-u,-w,-v)$. It leaves $P_b$ unchanged and sends
$R_b$ to $-R_b$. Whenever an odd graph jet is nonzero, these are distinct
oriented contact germs. Thus an orientation-sensitive inverse cannot
be deduced from the unmarked probability germs by simply removing an
evenness hypothesis. The mark in \eqref{eq:rawdata} supplies exactly
the odd information used by \eqref{eq:oddblock}. Even when all $R_b$
vanish, the recursion applies: in conjunction with the probability
germs it forces the successive odd jets to vanish.
\end{remark}
